%% Manuscript for AEU - International Journal of Electronics and Communications
%% Elsevier CAS template (cas-dc, double column), v2.4
%% Compile: pdflatex main; bibtex main; pdflatex main; pdflatex main
\documentclass[a4paper,fleqn]{cas-dc}

\usepackage[numbers,sort&compress]{natbib}
\usepackage{amsmath}
\usepackage{multirow}
\usepackage[siunitx]{circuitikz}
\usetikzlibrary{calc}
\usepackage{siunitx}
\sisetup{range-phrase=--, range-units=single}
\usepackage{xcolor}

% placeholder marker for items the authors must complete
\newcommand{\TBD}[1]{\textcolor{red}{[#1]}}

\begin{document}
\let\WriteBookmarks\relax
\def\floatpagepagefraction{1}
\def\textpagefraction{.001}

\shorttitle{Floating-gate floating-body resonator neuron in SKY130}
\shortauthors{J.L. Valtierra S\'anchez de la Vega}

\title [mode = title]{A programmable resonator (Type~II) spiking neuron in a single floating-gate, floating-body nMOS transistor: post-layout evidence from the open SKY130 process}

%% TODO authors: replace with the real author list, order, e-mails and ORCIDs
\author[1]{Jos\'e Luis Valtierra S\'anchez de la Vega}[orcid=0000-0000-0000-0000]
\cormark[1]
\ead{author.one@uabc.edu.mx}
\credit{\TBD{Conceptualization, Methodology, Software, Writing -- original draft}}

\author[1]{\TBD{Author Two}}
\credit{\TBD{Investigation, Validation}}

\author[1]{\TBD{Author Three}}
\credit{\TBD{Supervision, Writing -- review \& editing}}

\affiliation[1]{organization={Facultad de Ingenier\'ia Mexicali, Universidad Aut\'onoma de Baja California},
            city={Mexicali},
            state={Baja California},
            country={Mexico}}

\cortext[1]{Corresponding author}

\begin{abstract}
Floating-body transistors can emulate neurons because impact ionization charges their body and produces abrupt, hysteretic switching, but published demonstrations rely on measured devices or on behavioural models with fitted parameters. We present a neuron whose firing device is a single 5-V nMOS transistor with both a floating body, isolated in a deep n-well and discharged through a bulk-control transistor, and a capacitively coupled floating gate. The cell was laid out in the open SKY130 process (\qtyproduct{36 x 28}{\micro\meter}, DRC-clean), extracted with Magic, and simulated in ngspice with only the foundry compact models. The extracted cell shows static hysteresis whose trigger voltage moves from \SIrange{5.38}{4.26}{V} as the stored floating-gate voltage rises from \SIrange{0.40}{0.55}{V}. Driven by a constant current into its \SI{208}{fF} membrane capacitance, it fires at \SIrange{11}{210}{kHz} with \SIrange{2.5}{4.6}{pJ} per spike, inside a current band that the stored charge shifts. Firing starts at a finite rate, and just below onset the membrane oscillates at a nearly fixed \SI{15}{kHz} with an amplitude that grows with input current: the cell is a Type~II (resonator) neuron. Control-gate pulses gate firing below rheobase and mode-lock it above, but do not grade the output rate. Mismatch leaves 24 of 25 Monte Carlo samples firing at nominal bias, and trimming the stored charge, with the bulk-control bias as a second knob, restores firing at all process corners and temperatures studied. These results show, from foundry models and layout parasitics alone, that a standard open-PDK transistor can serve as a trimmable resonator neuron.
\end{abstract}

\begin{highlights}
\item One 5-V nMOS with a floating gate and a floating body fires as a spiking neuron
\item Evidence from SKY130 foundry models and extracted layout only, no macromodel
\item Stored floating-gate charge programs the trigger voltage from 5.38 V to 4.26 V
\item Fires at 11--210 kHz with 2.5--4.6 pJ per spike inside a programmable current band
\item Type II excitability: finite-rate onset, subthreshold oscillation, input mode locking
\end{highlights}

\begin{keywords}
Neuromorphic circuits \sep Floating-body MOSFET \sep Floating gate \sep Spiking neuron \sep Impact ionization \sep SKY130 \sep Post-layout simulation
\end{keywords}

\maketitle

%% ------------------------------------------------------------------
\section{Introduction}
\label{sec:intro}

Hardware spiking neural networks need large numbers of neurons, and the area of each neuron limits how many fit on a die. Conventional CMOS integrate-and-fire neurons use a membrane capacitor, a comparator or positive-feedback stage, and a reset path, which typically takes tens of transistors~\cite{indiveri2011,mead1989}. State-of-the-art neuromorphic processors use at least 18 transistors per neuron~\cite{pazos2025}. Reducing the neuron to one or two devices would raise density substantially.

Floating-body transistors are an attractive route. When the body of an nMOS transistor is not tied to a fixed potential, holes generated by impact ionization near the drain accumulate in the body, lower the threshold voltage through the body effect, and increase the drain current, which in turn generates more holes. This positive feedback produces a steep, hysteretic current jump known from silicon-on-insulator devices as the kink effect~\cite{boudou1987,park2018kink}. It has been used to build leaky integrate-and-fire neurons in partially depleted SOI~\cite{dutta2017,chavan2020} and single-transistor neurons co-integrated with synapses~\cite{han2021,woo2020}. Recently, Pazos~\textit{et al.} showed that a standard bulk nMOS transistor operated near punch-through, with its body connected to ground through a second transistor that sets the effective bulk resistance, reproduces neural and synaptic behaviours with low variability~\cite{pazos2025}.

Two gaps remain for circuit designers. First, these demonstrations are either measurements of specific devices or simulations with behavioural models whose parasitic-bipolar and avalanche parameters are fitted to those measurements. A designer working in an open process cannot easily tell whether the effect is available from the process design kit (PDK) itself. Second, floating-body neurons are sensitive to the body resistance and threshold voltage, so their firing threshold varies from die to die~\cite{pazos2025}. Floating-gate devices solve the analogous problem in analog circuits by storing a programmable, non-volatile offset on each transistor~\cite{diorio1996,bhattacharyya2025}, and floating-gate cells have been demonstrated in the open SkyWater 130-nm process~\cite{hasler2021woset,chen2024fg}.

This paper combines the two mechanisms in one transistor and asks a narrow question: \emph{does the open SKY130 PDK, with its foundry compact models and extracted layout parasitics alone, predict a working, programmable floating-body neuron?} We make five contributions:
\begin{enumerate}
\item A two-transistor neuron cell whose firing device is a 5-V nMOS transistor with both a floating gate and a floating body in a deep n-well, with a bulk-control transistor in the substrate (Section~\ref{sec:cell}).
\item A DRC-clean SKY130 layout and a parasitic-extracted netlist simulated only with the foundry models: no behavioural macromodel, no added bipolar device and no fitted parameter (Sections~\ref{sec:layout}--\ref{sec:method}).
\item Post-layout evidence of static hysteresis that does not depend on sweep rate, autonomous firing within a current band with depolarization block, and programming of the threshold and firing band by the stored floating-gate charge (Section~\ref{sec:results}).
\item A characterization of the excitability class: finite-frequency onset, a sustained subthreshold oscillation near onset, and mode locking to control-gate pulse trains, which identify the cell as a Type~II (resonator) neuron (Section~\ref{sec:typeii}).
\item A robustness study over process corners, temperature and Monte Carlo mismatch, showing that a per-cell floating-gate trim recovers firing where a fixed bias fails (Section~\ref{sec:robust}).
\end{enumerate}

%% ------------------------------------------------------------------
\section{Cell concept and operating principle}
\label{sec:cell}

\begin{figure}[t]
\centering
\resizebox{\linewidth}{!}{% Schematic of the floating-gate floating-body (FGFB) neuron cell
\begin{circuitikz}[american, font=\small]
  % input current and membrane capacitor
  \draw (0,6.8) node[above]{$V_{\mathrm{sup}}$} to[short,o-] (0,6.5)
        to[I, l_=$I_{\mathrm{in}}$] (0,5.0) coordinate (mem);
  \draw (mem) to[C, l_=$C_{\mathrm{mem}}$] (0,2.4) node[ground]{};
  \fill (mem) circle (1.6pt);
  \node[above right] at (mem) {$V_{\mathrm{mem}}$};
  % FGFB transistor M1
  \node[nmos, bulk] (M1) at (6.4,3.2) {};
  \draw (mem) -| (M1.D);
  \draw (M1.S) node[ground]{};
  \node[anchor=west] at ($(M1.D)+(0.2,-0.25)$) {$M_1$};
  % floating-gate network
  \draw (M1.G) -- ++(-1.4,0) coordinate (fg);
  \fill (fg) circle (1.6pt);
  \node[above] at (fg) {FG};
  \draw (fg) to[C, l=$C_{\mathrm{fix}}$] (fg |- 0,1.6) node[ground]{};
  \draw (fg) -- ++(-1.3,0) coordinate (ca);
  \draw (ca) to[C, l=$C_{A}$, -o] (ca |- 0,1.6) node[below]{$V_{A}$};
  \draw (ca) -- ++(-1.3,0) coordinate (cb);
  \draw (cb) to[C, l=$C_{B}$, -o] (cb |- 0,1.6) node[below]{$V_{B}$};
  \fill (ca) circle (1.6pt);
  \node[font=\footnotesize, align=center, anchor=north] at ($(fg)+(0,-2.15)$) {stored charge\\$\Rightarrow V_Q$};
  % body and bulk-control transistor M2
  \draw (M1.bulk) -- ++(1.1,0) coordinate (b);
  \fill (b) circle (1.6pt);
  \node[above] at (b) {body};
  \draw (b) -- ++(0,-0.5) node[nmos, anchor=D, xscale=-1] (M2) {};
  \draw (M2.S) node[ground]{};
  \draw (M2.G) -- ++(0.3,0) node[right]{$V_{G2}$};
  \node[anchor=south] at ($(M2.G)+(0.45,0.05)$) {$M_2$};
  % isolation box
  \draw[dashed, rounded corners] ($(M1.G)+(-0.15,0.95)$) rectangle ($(M1.bulk)+(0.45,-0.95)$);
  \node[font=\scriptsize, align=left, anchor=south west] at ($(M1.G)+(-0.15,0.95)$) {isolated p-well (deep n-well at 5\,V)};
\end{circuitikz}
}
\caption{Floating-gate floating-body (FGFB) neuron cell. $M_1$ is a SKY130 5-V nMOS transistor in an isolated p-well (deep n-well biased at \SI{5}{V}); its floating gate (FG) is coupled to two control inputs $V_A$, $V_B$ and to ground through MIM capacitors, and its stored charge is represented by the equivalent voltage $V_Q$. $M_2$, a second 5-V nMOS in the substrate, discharges the body with a gate bias $V_{G2}$. The input current $I_{\mathrm{in}}$ charges the membrane capacitance $C_{\mathrm{mem}}$ at the drain of $M_1$.}
\label{fig:schematic}
\end{figure}

Fig.~\ref{fig:schematic} shows the cell. The firing transistor $M_1$ sits in a p-well isolated from the substrate by a deep n-well, so its body is a separate node. The body is discharged only through the junctions of $M_1$ and through the bulk-control transistor $M_2$, whose gate voltage $V_{G2}$ sets the effective body resistance, as in~\cite{pazos2025}. The gate of $M_1$ is a floating node with no DC path; its voltage is
\begin{equation}
V_{FG} = V_Q + \beta_A V_A + \beta_B V_B + \beta_D V_{\mathrm{mem}},
\label{eq:vfg}
\end{equation}
with $\beta_i = C_i/C_{\Sigma}$, where $V_Q = Q_{FG}/C_\Sigma$ is the equivalent voltage of the stored charge, $C_\Sigma$ is the total floating-gate capacitance, and $\beta_D$ accounts for the gate--drain overlap and wiring coupling.

With the drain at several volts, impact ionization near the drain injects a hole current $I_{ii}$ into the body. The SKY130 compact models describe it with the BSIM4 substrate-current expression~\cite{bsim4}
\begin{equation}
\begin{split}
I_{ii} = {} & \frac{\alpha_0 + \alpha_1 L}{L}\,\left(V_{DS}-V_{DS,\mathrm{eff}}\right)\\
& \times \exp\!\left(\frac{-\beta_0}{V_{DS}-V_{DS,\mathrm{eff}}}\right) I_{DS},
\end{split}
\label{eq:iii}
\end{equation}
with foundry-extracted $\alpha_0$, $\alpha_1$ and $\beta_0$. The body voltage $V_B$ and the membrane voltage then obey
\begin{align}
C_B \frac{dV_B}{dt} &= I_{ii}(V_{\mathrm{mem}}, V_B, V_{FG}) \nonumber\\
 &\quad - I_{SB}(V_B) - I_{M2}(V_B, V_{G2}), \label{eq:body}\\
C_{\mathrm{mem}} \frac{dV_{\mathrm{mem}}}{dt} &= I_{\mathrm{in}} - I_{D}(V_{\mathrm{mem}}, V_B, V_{FG}), \label{eq:mem}
\end{align}
where $I_{SB}$ is the source--body junction current and $C_B$ the body capacitance. Because a higher $V_B$ lowers the threshold voltage and raises $I_D$ and hence $I_{ii}$, Eq.~\eqref{eq:body} contains positive feedback. When the loop gain exceeds unity the body switches to a high state, clamped by the forward-biased source junction; this produces a jump in drain current at a trigger voltage and a return at a lower release voltage. All of this behaviour is contained in the standard BSIM4 equations; no lateral bipolar transistor is added. The stored charge $V_Q$ sets the channel current that feeds Eq.~\eqref{eq:iii} and therefore the trigger voltage, while the control inputs shift $V_{FG}$ dynamically through Eq.~\eqref{eq:vfg}.

When $C_{\mathrm{mem}}$ is charged by $I_{\mathrm{in}}$, the cell is a two-variable system, with the body as the fast variable and the membrane as the slow one. If the steady state falls on the unstable middle branch of the hysteretic characteristic, the cell oscillates: $V_{\mathrm{mem}}$ rises to the trigger voltage, the body switches, the drain current discharges $C_{\mathrm{mem}}$ to the release voltage, the body relaxes through $M_2$, and the cycle repeats.

%% ------------------------------------------------------------------
\section{Layout and parasitic extraction}
\label{sec:layout}

\begin{figure}[t]
\centering
\includegraphics[width=\linewidth]{figs/fig_layout.pdf}
\caption{SKY130 layout of the FGFB neuron cell (\qtyproduct{36.3 x 28.1}{\micro\meter}, DRC-clean in Magic). The floating gate is routed on metal~4 to the top plates of the MIM capacitors; the membrane node is routed on metal~2 to $C_{\mathrm{mem}}$.}
\label{fig:layout}
\end{figure}

The cell was drawn in Magic~\cite{magic} with the SKY130A technology files from open\_pdks~\cite{openpdks} (Fig.~\ref{fig:layout}). Both transistors are \texttt{sky130\_fd\_pr\_\_\allowbreak{}nfet\_g5v0d10v5} devices with $W/L = \SI{1}{\micro\meter}/\SI{0.5}{\micro\meter}$, generated by the PDK device generator. $M_1$ and its p-type guard ring, which forms the body contact, are enclosed in a deep n-well surrounded by an n-well ring tied to \SI{5}{V}. The body contact is routed on metal~1 to the drain of $M_2$, which sits in the substrate.

The coupling capacitors are \texttt{sky130\_fd\_pr\_\_\allowbreak{}cap\_mim\_m3\_1} devices. The smallest MIM capacitor (\qtyproduct{2 x 2}{\micro\meter}, \SI{8.4}{fF}) sets the control-gate capacitances $C_A = C_B$, and $C_{\mathrm{fix}}$ (\qtyproduct{9.6 x 9.6}{\micro\meter}, \SI{192}{fF}) sets the coupling ratio to $\beta_{A,B}\approx 0.04$. The floating gate connects only to the polysilicon gate and to the MIM top plates, which minimizes its parasitic capacitance to the substrate. $C_{\mathrm{mem}}$ is a \qtyproduct{10 x 10}{\micro\meter} MIM capacitor (\SI{208}{fF}). The cell passes full DRC, and its extracted connectivity matches the schematic.

Parasitic capacitances were extracted with Magic (\texttt{extract all}, \texttt{ext2spice cthresh 0}). Table~\ref{tab:pex} summarizes the values that matter for the neuron. The p-well to deep-n-well junction dominates the body capacitance, and routing over the isolated well adds \SI{1.6}{fF} between the floating gate and the body. Direct wiring coupling between drain and floating gate is \SI{0.2}{fF}, in addition to the overlap capacitance contained in the device model. In transient simulation, a \SI{1.8}{V} step on one control input moves the floating gate by \SIrange{77}{82}{mV}, i.e.\ an effective $\beta_A \approx 0.044$.

\begin{table}[t]
\centering
\caption{Extracted capacitances relevant to the neuron (Magic, SKY130A). Device-intrinsic junction and overlap capacitances are inside the BSIM4 models and are not listed.}
\label{tab:pex}
\begin{tabular}{@{}lll@{}}
\toprule
Node & Component & Value (fF) \\
\midrule
\multirow{3}{*}{Body} & p-well to deep n-well (5-V ring) & 12.6 \\
 & wiring to $V_{\mathrm{mem}}$ / FG & 2.2 / 1.6 \\
 & total extracted & $\approx$16.8 \\
\midrule
\multirow{3}{*}{Floating gate} & MIM: $C_{\mathrm{fix}}$, $C_A$, $C_B$ & 192, 8.4, 8.4 \\
 & wiring and plate parasitics & $\approx$3 \\
 & of which to $V_{\mathrm{mem}}$ (drain) & 0.2 \\
\midrule
\multirow{2}{*}{Membrane} & MIM $C_{\mathrm{mem}}$ & 208 \\
 & metal fringe and wiring & $\approx$12.5 \\
\bottomrule
\end{tabular}
\end{table}

%% ------------------------------------------------------------------
\section{Simulation methodology}
\label{sec:method}

All results use the extracted netlist and ngspice~42~\cite{ngspice} with the SKY130 compact models built by open\_pdks~\cite{sky130pdk,openpdks}, typical corner and \SI{27}{\celsius} unless stated otherwise. The netlist is simulated exactly as extracted; no element is added to the transistors. Two test-bench elements stand in for things that are outside the PDK: a \SI{e18}{\ohm} resistor from the floating gate to a source $V_Q$, which sets the stored charge in the DC operating point and has a time constant of order \SI{e5}{s}, and an ideal current source for $I_{\mathrm{in}}$. Floating-gate programming (tunnelling and injection) is not modelled; $V_Q$ is treated as a programmed value. Unless stated, $V_{G2} = \SI{0.5}{V}$.

Hysteresis was measured with triangular drain-voltage sweeps from \SIrange{0.05}{6}{V}. A sweep that is fast compared with the body time constant also produces a loop, purely from the lag of $V_B$ behind the drain voltage; such a loop shrinks in proportion to the sweep rate. We therefore classify a loop as static only if it is unchanged when the ramp time is increased by two decades. For transients, the maximum time step was \SI{4}{ns} on an \SI{80}{\micro\second} run. Halving and doubling it changed the firing rate by less than 0.01\,\%.

The SKY130 impact-ionization parameters in Eq.~\eqref{eq:iii} are extracted by the foundry with the body grounded. Using them with a floating body near the punch-through regime is an extrapolation that we cannot verify without measurement, and we treat it as the main limitation of this study (Section~\ref{sec:discussion}).

%% ------------------------------------------------------------------
\section{Post-layout results}
\label{sec:results}

\subsection{Static hysteresis programmed by stored charge}

\begin{figure*}[t]
\centering
\includegraphics[width=0.92\textwidth]{figs/fig_hysteresis.pdf}
\caption{Quasi-static $I_D$--$V_{\mathrm{mem}}$ characteristics of the extracted cell (stock SKY130 models, $V_{G2}=\SI{0.5}{V}$). (a) The stored floating-gate voltage $V_Q$ moves the hysteresis loop (\SI{2}{ms} ramps). (b) At $V_Q=\SI{0.5}{V}$ the loop is essentially unchanged between a \SI{200}{\micro\second} and a \SI{20}{ms} ramp, which identifies it as static bistability rather than sweep-rate lag.}
\label{fig:hyst}
\end{figure*}

Fig.~\ref{fig:hyst}(a) shows the extracted cell swept slowly at four values of stored charge. Each curve jumps from tens of nanoamperes to the microampere range at a trigger voltage and returns at a lower release voltage. Table~\ref{tab:hyst} lists the extracted trigger and release voltages. Raising $V_Q$ from \SIrange{0.40}{0.55}{V} lowers the trigger from \SIrange{5.38}{4.26}{V}, a sensitivity of about \SI{7.5}{V/V}, and narrows the loop from \SIrange{0.92}{0.53}{V}. Fig.~\ref{fig:hyst}(b) compares a \SI{200}{\micro\second} and a \SI{20}{ms} ramp. The loop area falls only from \SIrange{0.53}{0.34}{\micro\ampere\volt}, and a \SI{2}{ms} ramp already gives \SI{0.37}{\micro\ampere\volt}, so the loop converges to a finite static value. By contrast, a loop dominated by body lag shrinks roughly tenfold per decade of ramp time.

The bulk-control bias adds a second degree of freedom (Table~\ref{tab:hyst}, lower part). At fixed $V_Q$, raising $V_{G2}$ raises the trigger voltage and the current ratio between the high and low branches. Over the explored range, static loops exist for $V_Q$ between \SIrange{0.3}{0.7}{V} and $V_{G2}$ between \SIrange{0.3}{0.6}{V}, with high-to-low current ratios of 5 to 39.

\begin{table*}[t]
\centering
\caption{Static hysteresis of the extracted cell. Top: trigger and release voltage versus stored charge at $V_{G2}=\SI{0.5}{V}$. Bottom: trigger voltage (V) and maximum high/low current ratio versus $V_Q$ and $V_{G2}$ (\SI{2}{ms} ramps; ``--'' means no loop below \SI{5.9}{V}).}
\label{tab:hyst}
\begin{tabular}{@{}lcccc@{}}
\toprule
$V_Q$ (V) & 0.40 & 0.45 & 0.50 & 0.55 \\
\midrule
Trigger (V) & 5.38 & 4.95 & 4.58 & 4.26 \\
Release (V) & 4.45 & 4.18 & 3.94 & 3.73 \\
Window (V)  & 0.92 & 0.77 & 0.64 & 0.53 \\
\bottomrule
\end{tabular}

\vspace{0.6em}
\begin{tabular}{@{}lcccc@{}}
\toprule
$V_Q$ \textbackslash{} $V_{G2}$ & 0.3\,V & 0.4\,V & 0.5\,V & 0.6\,V \\
\midrule
0.3\,V & 4.73 / 13 & 5.42 / 26 & -- & -- \\
0.4\,V & 4.12 / 14 & 4.62 / 23 & 5.38 / 39 & -- \\
0.5\,V & 3.64 / 11 & 4.01 / 17 & 4.58 / 25 & 5.37 / 38 \\
0.6\,V & 3.26 / 8  & 3.55 / 11 & 3.98 / 15 & 4.58 / 20 \\
0.7\,V & 2.95 / 6  & 3.18 / 7  & 3.52 / 8  & 3.99 / 10 \\
\bottomrule
\end{tabular}
\end{table*}

\subsection{Spiking operation and firing band}

\begin{figure*}[t]
\centering
\includegraphics[width=0.92\textwidth]{figs/fig_neuron.pdf}
\caption{Neuron operation of the extracted cell. (a) Membrane and body voltages for $V_Q=\SI{0.45}{V}$ and $I_{\mathrm{in}}=\SI{50}{nA}$ (\SI{85.4}{kHz}). (b) Firing rate versus input current for four values of stored charge; crosses mark non-firing points (leak-limited at low current, depolarization block at high current).}
\label{fig:neuron}
\end{figure*}

With a constant input current into the layout membrane capacitance, the cell fires without any external clock or reset (Fig.~\ref{fig:neuron}(a)). At $V_Q = \SI{0.45}{V}$ and $I_{\mathrm{in}} = \SI{50}{nA}$, $V_{\mathrm{mem}}$ ramps from \SIrange{3.9}{5.2}{V}, the body pulses to about \SI{0.55}{V}, and the cell fires at \SI{85.4}{kHz}. The ramp is set by $I_{\mathrm{in}}/C_{\mathrm{mem}}$ and the fast fall by the switching of the body, which matches the slow--fast picture of Section~\ref{sec:cell}.

Firing occurs only inside a band of input currents (Fig.~\ref{fig:neuron}(b)). Below the band, the subthreshold leakage of $M_1$ at the trigger voltage exceeds $I_{\mathrm{in}}$, and $V_{\mathrm{mem}}$ settles below the trigger. Above it, the steady state moves onto a stable high-current branch and the cell stops firing at an elevated membrane voltage, analogous to depolarization block in biological neurons~\cite{izhikevich2007}. Inside the band, the rate grows almost linearly with $I_{\mathrm{in}}$, by about \SI{1.7}{kHz/nA} at $V_Q = \SI{0.45}{V}$. The stored charge shifts the band to higher currents: approximately \SIrange{10}{50}{nA} at $V_Q=\SI{0.40}{V}$, \SIrange{25}{100}{nA} at \SI{0.45}{V} and \SIrange{50}{150}{nA} at \SI{0.50}{V}. At \SI{0.55}{V}, the cell did not fire up to \SI{450}{nA}, although its static loop persists. This shows that hysteresis is necessary but not sufficient: the oscillation also requires the body to be faster than the membrane.

The energy supplied by the input per spike, $E = I_{\mathrm{in}}\langle V_{\mathrm{mem}}\rangle/f$, is \SIrange{2.5}{3.1}{pJ} across the firing band and rises to \SI{4.6}{pJ} at the lowest rate (\SI{10.9}{kHz}, \SI{10}{nA}). This figure excludes the static current of the deep-n-well bias and the gate bias of $M_2$, which draw only junction leakage.

\subsection{Control-gate inputs}

\begin{figure*}[t]
\centering
\includegraphics[width=0.92\textwidth]{figs/fig_modulation.pdf}
\caption{Effect of control-gate inputs at $V_Q=\SI{0.45}{V}$ and $I_{\mathrm{in}}=\SI{50}{nA}$. Raising $V_A$ (green, floating gate $+$\SI{80}{mV}) silences the cell; lowering $V_B$ (red, floating gate $-$\SI{80}{mV}) keeps it firing with a higher threshold.}
\label{fig:mod}
\end{figure*}

The two control gates allow inputs to act directly on the firing transistor (Fig.~\ref{fig:mod}). A \SI{1.8}{V} step on $V_A$ raises the floating gate by about \SI{80}{mV}. The extra leakage then exceeds $I_{\mathrm{in}}$, and the cell stops firing and settles at about \SI{4.4}{V}. Lowering $V_B$ by \SI{1.8}{V} lowers the floating gate by the same amount. The cell keeps firing with a higher peak voltage (\SI{5.6}{V} instead of \SI{5.2}{V}) and a slightly higher rate (\SI{94}{kHz} instead of \SI{85}{kHz}). A positive control-gate input therefore acts as shunting inhibition, because at fixed input current the increase in leakage dominates over the decrease in trigger voltage. This polarity is the opposite of what a simple threshold-shift picture predicts, and it is a direct consequence of the foundry model.

\subsection{Excitability class and response to gate pulses}
\label{sec:typeii}

\begin{figure*}[t]
\centering
\includegraphics[width=0.92\textwidth]{figs/fig_typeII.pdf}
\caption{Excitability class of the extracted cell. (a) Firing rate near onset; open circles are non-firing points (membrane swing below \SI{0.4}{V}) and dotted lines mark the jump from rest to a finite rate. (b) Membrane voltage just above onset at $V_Q=\SI{0.40}{V}$, $V_{G2}=\SI{0.4}{V}$: the oscillation starts with small amplitude at about \SI{15}{kHz} and grows with $I_{\mathrm{in}}$ while its frequency hardly changes. (c) Output rate for a \SI{100}{kHz} train of active-low pulses on $V_A$ at $V_Q=\SI{0.45}{V}$ and $I_{\mathrm{in}}=\SI{20}{nA}$, a bias that is silent without input; the rate moves between the $f_{in}/4$ and $f_{in}/3$ locking plateaus. (d) Output rate with the same \SI{2}{\micro\second} pulses above rheobase, $I_{\mathrm{in}}=\SI{40}{nA}$.}
\label{fig:typeii}
\end{figure*}

How a neuron begins to fire determines how it encodes its input~\cite{izhikevich2007}. A Type~I (integrator) neuron starts at an arbitrarily low rate, so its rate grades with input strength. A Type~II (resonator) neuron starts at a finite rate through a Hopf-like bifurcation, often preceded by small subthreshold oscillations, and responds preferentially to inputs near its own frequency. We characterized the extracted cell near onset with \SIrange{2}{3}{ms} transients and counted only membrane swings above \SI{0.4}{V} as spikes.

\textbf{Finite-frequency onset.} Firing never starts at a low rate (Fig.~\ref{fig:typeii}(a)). At $V_Q=\SI{0.40}{V}$ the cell is silent up to \SI{8}{nA} and fires at \SI{10.9}{kHz} at \SI{10}{nA}; at $V_Q=\SI{0.38}{V}$ it is silent at \SI{5}{nA} and fires at \SI{9.2}{kHz} at \SI{8}{nA}. With $V_{G2}=\SI{0.4}{V}$ the first spikes appear at \SI{15.2}{kHz}. At $V_{G2}=\SI{0.3}{V}$ and \SI{0.6}{V} the cell did not fire at $V_Q=\SI{0.40}{V}$ for \SIrange{4}{18}{nA}, so no bulk bias in this range produced an integrator-like onset.

\textbf{Subthreshold oscillation.} At $V_{G2}=\SI{0.4}{V}$ the transition was resolved more finely (Fig.~\ref{fig:typeii}(b)). At \SI{7}{nA} perturbations decay. From \SIrange{7.5}{9}{nA} the membrane oscillates continuously with peak-to-peak amplitudes of 0.14, 0.31, 0.45 and \SI{0.58}{V} at a nearly constant \SIrange{15.1}{15.4}{kHz}, and at \SI{12}{nA} the swing reaches \SI{1.14}{V} at \SI{17.8}{kHz}. An amplitude that grows from zero at a fixed, finite frequency is the behaviour expected near a supercritical Hopf bifurcation; the amplitudes follow roughly $\sqrt{I_{\mathrm{in}}-I_H}$ with $I_H\approx\SI{7.4}{nA}$. The \SI{8}{nA} amplitude changed by less than \SI{1}{mV} when the maximum time step was reduced from 20 to \SI{5}{ns}, so the oscillation is not a numerical artefact.

\textbf{Gate pulses below rheobase.} We tested the control gates as a possible synaptic input. A falling edge on $V_A$ lowers the floating gate by about \SI{80}{mV} and transiently raises the trigger voltage; when the pulse ends, the membrane is left above the restored trigger and the cell fires. At $V_Q=\SI{0.45}{V}$ and $I_{\mathrm{in}}=\SI{20}{nA}$, which is silent without input, a train of \SI{2}{\micro\second} pulses evokes firing that stops when the train stops. The output rate, however, does not grade with input strength (Fig.~\ref{fig:typeii}(c)). Pulse amplitudes from \SIrange{0.6}{1.8}{V} all give \SI{25.0}{kHz}, one spike per four input pulses; only \SI{0.3}{V} gives a different value, \SI{23.1}{kHz}. Pulse width moves the rate in steps, from \SI{21.0}{kHz} at \SI{0.5}{\micro\second} to the $f_{in}/4$ plateau at \SIrange{1}{2}{\micro\second} and the $f_{in}/3$ plateau at \SIrange{4}{8}{\micro\second}. The stored charge acts as a switch rather than as a scale factor: the same input is passed at $V_Q\le\SI{0.45}{V}$ and blocked from about \SI{0.455}{V} upward. Close to onset ($V_Q=\SI{0.38}{V}$) the result is the same: at \SI{6}{nA}, trains at \SIrange{5}{40}{kHz} never evoke firing, and at \SI{8}{nA} they pull the intrinsic \SI{9.2}{kHz} rhythm to \SIrange{9.3}{10.0}{kHz}.

\textbf{Mode locking above rheobase.} With the cell already firing at \SI{40}{nA}, pulse trains move the rate onto rational fractions of the input frequency (Fig.~\ref{fig:typeii}(d)). A \SI{200}{kHz} train locks all three stored-charge values to \SI{66.7}{kHz} ($f_{in}/3$), and a \SI{50}{kHz} train pulls the $V_Q=\SI{0.48}{V}$ cell from \SI{58.2}{kHz} to exactly \SI{50.0}{kHz}. These plateaus, and the steps between them in Fig.~\ref{fig:typeii}(c), are the Arnold tongues of a driven nonlinear oscillator.

Taken together, the cell behaves as a resonator: it fires from a finite frequency, oscillates below threshold, and synchronizes to periodic input. The control gates therefore implement gating and entrainment, which are useful for coincidence detection and frequency-selective responses~\cite{izhikevich2007}, but not a graded synaptic weight, because the output rate is quantized by locking rather than proportional to the input. We searched for a graded regime over stored charge, bulk bias (\SIrange{0.3}{0.6}{V}), input current (from below onset to mid-band), and pulse amplitude, width and rate, and found none in the extracted cell.

%% ------------------------------------------------------------------
\section{Robustness and floating-gate trimming}
\label{sec:robust}

\begin{figure*}[t]
\centering
\includegraphics[width=0.92\textwidth]{figs/fig_robust.pdf}
\caption{Robustness of the extracted cell at $I_{\mathrm{in}}=\SI{50}{nA}$. (a) Firing points versus stored charge for the five process corners (\SI{27}{\celsius}) and three temperatures (typical corner); colour gives the firing rate, diamonds mark points that fire after raising $V_{G2}$ to \SIrange{0.6}{0.7}{V}. (b) Twenty-five Monte Carlo mismatch samples (typical corner, \SI{27}{\celsius}, $V_{G2}=\SI{0.5}{V}$).}
\label{fig:robust}
\end{figure*}

Floating-body neurons are sensitive to the threshold voltage and the body resistance. We evaluated this with the SKY130 corner models and with its Monte Carlo mismatch models, which apply Gaussian variations to the threshold voltage, oxide thickness and subthreshold parameters of each device instance. Each condition was simulated at $I_{\mathrm{in}} = \SI{50}{nA}$ over a grid of stored-charge values.

\textbf{Mismatch.} Device mismatch is small compared with the width of the firing window (Fig.~\ref{fig:robust}(b)). Of 25 Monte Carlo samples, 24 fire at the nominal $V_Q = \SI{0.45}{V}$, at $77.8 \pm 7.6$\,kHz (range \SIrange{66.6}{93.1}{kHz}), and the remaining sample fires after a trim to an adjacent grid point. Each sample fires over one to four of the five grid points (\SIrange{0.35}{0.55}{V}).

\textbf{Process corners and temperature.} Global variations move the window much further (Fig.~\ref{fig:robust}(a) and Table~\ref{tab:trim}). At $V_Q = \SI{0.45}{V}$ and $V_{G2}=\SI{0.5}{V}$, only the typical corner at \SI{27}{\celsius} fires. For the fast-nMOS corners the window moves to lower stored charge, to \SIrange{0.25}{0.35}{V} for ff and \SIrange{0.15}{0.25}{V} for sf, and at \SI{85}{\celsius} it moves to \SIrange{0.25}{0.35}{V}. In all three cases a stored-charge trim alone restores firing. For the slow-nMOS corners (ss, fs) and at \SI{-40}{\celsius}, the static loop persists but the oscillation is not sustained anywhere on the $V_Q$ grid at $V_{G2} = \SI{0.5}{V}$; the body is then too slow relative to the membrane. Raising $V_{G2}$ to \SIrange{0.6}{0.7}{V} lowers the body resistance and restores firing at \SIrange{81}{93}{kHz} (Table~\ref{tab:trim}). The ss corner also fires without changing $V_{G2}$ if the input current is lowered to \SI{20}{nA} (\SI{26.8}{kHz} at $V_Q = \SI{0.55}{V}$), consistent with the slow--fast condition of Section~\ref{sec:cell}.

\begin{table}[t]
\centering
\caption{Trim points that restore firing at $I_{\mathrm{in}} = \SI{50}{nA}$ (post-layout, stock SKY130 models). Only the typical corner at \SI{27}{\celsius} fires at the nominal bias.}
\label{tab:trim}
\begin{tabular}{@{}lccc@{}}
\toprule
Condition & $V_Q$ (V) & $V_{G2}$ (V) & Rate (kHz) \\
\midrule
tt, \SI{27}{\celsius} (nominal) & 0.40--0.50 & 0.5 & 77.7--85.6 \\
ff & 0.25--0.35 & 0.5 & 77.4--93.2 \\
sf & 0.15--0.25 & 0.5 & 69.4--76.1 \\
ss & 0.55 & 0.6--0.7 & 81.3--87.9 \\
fs & 0.55 & 0.7 & 82.5 \\
tt, \SI{85}{\celsius} & 0.25--0.35 & 0.5 & 97.7--116.4 \\
tt, \SI{-40}{\celsius} & 0.45 & 0.6 & 93.1 \\
\bottomrule
\end{tabular}
\end{table}

Because the firing rate inside the band is set mainly by $I_{\mathrm{in}}$ and $C_{\mathrm{mem}}$, a trimmed cell recovers not only firing but a similar rate: the trimmed conditions fire at \SIrange{69}{116}{kHz} against \SI{85}{kHz} nominal. The required stored-charge range, \SIrange{0.15}{0.55}{V}, corresponds to a few hundred millivolts of floating-gate offset, which is within the range used by floating-gate analog circuits~\cite{diorio1996,hasler2021woset}. In a fabricated cell, $V_Q$ would be programmed once per neuron and $V_{G2}$ shared by a block of neurons, which is the main motivation for placing the floating gate on the firing transistor itself.

%% ------------------------------------------------------------------
\section{Comparison and discussion}
\label{sec:discussion}

\begin{table*}[t]
\centering
\caption{Comparison with transistor-level spiking neurons.}
\label{tab:compare}
{%
\begin{tabular}{@{}lccc@{}}
\toprule
 & Pazos \textit{et al.}~\cite{pazos2025} & Bhattacharyya \& Hasler~\cite{bhattacharyya2025} & This work \\
\midrule
Evidence & measured & measured & post-layout sim. \\
Process & bulk CMOS, 180/500 nm & 350 nm FPAA & SKY130, 130 nm \\
Transistors & 2 & 6 & 2 (+ 4 MIM) \\
Mechanism & floating body & FG Schmitt trigger & floating body + FG \\
Programmability & $V_{G2}$ & FG charge & FG charge + $V_{G2}$ \\
Firing rate & --- & 2\,Hz--20\,kHz & 11--210\,kHz \\
Energy / power & 415\,pJ/\textmu m at 12.6\,\textmu s & 1.1--2.6\,\textmu W & 2.5--4.6\,pJ/spike \\
Area & --- & 0.0025\,mm$^2$ & 0.0010\,mm$^2$ (cell) \\
Supply / drive & 3.6--4.5\,V pulses & --- & 4--5.6\,V on $V_{\mathrm{mem}}$ \\
\bottomrule
\end{tabular}}
\end{table*}

Table~\ref{tab:compare} places the cell among transistor-level neurons. Its distinguishing features are the combination of floating-body firing and floating-gate programming in one device, and the fact that the evidence comes entirely from an open PDK's models and extracted layout. The area of \SI{0.0010}{mm^2} is dominated by the MIM capacitors; the two transistors with their isolation occupy roughly \qtyproduct{9 x 9}{\micro\meter}. Smaller coupling capacitors (e.g.\ metal-fringe capacitors) or a smaller membrane capacitor would shrink the cell and raise the firing rate.

\textbf{Limitations.} First, and most important, the impact-ionization parameters of the SKY130 models are extracted with a grounded body under normal bias. The floating-body regime used here is an extrapolation of a calibrated model, not a calibrated prediction, and the absolute trigger voltages and currents should be confirmed by measurement or device simulation. Second, the models contain no lateral bipolar transistor. Measured floating-body devices show larger snapback currents than our cell~\cite{pazos2025}, so the BSIM4-only behaviour is likely a conservative lower bound on the feedback. Third, the membrane reaches \SIrange{5.2}{5.6}{V}, close to the nominal rating of the 5-V device, so long-term reliability would need assessment. Fourth, floating-gate programming circuitry and charge retention are not modelled; $V_Q$ is an ideal programmed value. Finally, the control gates do not provide a graded synaptic weight: a sustained positive input acts as shunting inhibition, and pulsed negative inputs gate or entrain firing with an output rate quantized by mode locking (Section~\ref{sec:typeii}). Weighted synaptic input must therefore still be supplied as current into $C_{\mathrm{mem}}$, for example from separate floating-gate synapse transistors.

%% ------------------------------------------------------------------
\section{Conclusion}
\label{sec:conclusion}

We have shown, using only the open SKY130 compact models and parasitics extracted from a DRC-clean layout, that a 5-V nMOS transistor with a floating body and a floating gate operates as a programmable spiking neuron. The extracted cell exhibits static hysteresis whose trigger voltage is programmed by the stored floating-gate charge (\SIrange{5.38}{4.26}{V} for \SIrange{0.40}{0.55}{V}). With its own \SI{208}{fF} membrane capacitor, it fires at \SIrange{11}{210}{kHz} with \SIrange{2.5}{4.6}{pJ} per spike inside an input-current band that the stored charge shifts, and it shows depolarization block and shunting inhibition through its control gates. Its onset is that of a Type~II resonator: firing starts at \SIrange{9}{15}{kHz}, a sustained subthreshold oscillation appears just below, and control-gate pulse trains gate and mode-lock the output rather than weighting it. Mismatch alone leaves 24 of 25 samples firing at the nominal bias; process corners and temperature move the firing window by more than its width, but trimming the stored charge, with the bulk-control bias as a second knob, recovers firing in all eight conditions studied. Measurements of fabricated cells are the necessary next step to confirm the floating-body regime of the foundry models.

%% ------------------------------------------------------------------
\section*{Declaration of competing interest}
The authors declare that they have no known competing financial interests or personal relationships that could have appeared to influence the work reported in this paper.

\section*{Data availability}
The layout (GDS and Magic files), extracted netlists, test benches and analysis scripts will be made available in a public repository \TBD{(URL)}.

\section*{Declaration of generative AI and AI-assisted technologies in the writing process}
\TBD{During the preparation of this work the authors used [tool] in order to [purpose]. After using this tool, the authors reviewed and edited the content as needed and take full responsibility for the content of the publication.}


\printcredits

\bibliographystyle{elsarticle-num}
\bibliography{refs}

\end{document}
